\scanpage{172}\mspage{27}%
going to the very end\unsure{,} of this skin-deep resistance, one comes to generalize and deepen the notion of ``homotopy type'', and to bring yet more fresh water to the development of a good homotopical yoga. Here is what I am getting at.

Let us first speak of sheaves \inserted{(of sets, or of modules and so on)}, for simplicity, instead of stacks, \inserted{and \unsure{let us work|let us work} on the étale topos of a scheme $X$.} Locally constant sheaves - modulo a fairly innocuous additional finiteness condition - form the simplest of the \uline{constructible} sheaves, for whose definition they serve as a model. Assuming $X$ coherent ($=$ quasi-compact and quasi-separated), general constructible sheaves are those for which there exists a \inserted{finite} partition $X=\bigcup\limits_{i\in I}X_i$\raisebox{-1.7ex}[0pt][0pt]{\makebox[0pt][r]{\parbox{1.5mm}{\figure{\begin{tikzpicture}[x=1mm,y=1mm]\fill (0,0) circle[radius=.14];\end{tikzpicture}}}\hspace{.8em}}} of $X$ into ``cells'' \inserted{or ``strata''} $X_i$, locally closed and constructible, such that \deleted{\unsure{each|each}} the restriction of $F$ to each $X_i$ is locally constant \deleted{and} (plus a finiteness condition\ldots) Thus the category of constructible sheaves on $X$ . (which gives back the category of all sheaves by passing to a category of ind-objects\ldots) it-$\overline m$ \unsure{appears|appears} as an inductive limit of categories, associated with partitions more and
% Small point at whole-image (1161,1361), below/right of the union index: reproduced by the zero-width raised box; no mathematical role is assigned.
% The separate point after X on L20 is reproduced as a physical point, not asserted to end a sentence.
% archive sheet number: 172
